\documentclass[12pt, a0paper]{article}

% ===== 页面 (PPT 模板: 84cm × 120cm) =====
\usepackage[
    paperwidth=84cm,
    paperheight=120cm,
    margin=1.8cm,
    columnsep=1.8cm
]{geometry}

% ===== 中文 =====
\usepackage{ctex}

% ===== 双栏 =====
\usepackage{multicol}

% ===== 表格 =====
\usepackage[table]{xcolor}
\usepackage{colortbl}
\usepackage{multirow}
\usepackage{booktabs}
\usepackage{tabularx}

% ===== 图形 =====
\usepackage{tikz}
\usetikzlibrary{shapes.geometric, backgrounds, positioning}

% ===== 其他 =====
\usepackage{enumitem}
\usepackage{amsmath}
\usepackage{amssymb}
\usepackage{fontawesome5}

% ===== PPT 模板配色 =====
\definecolor{myteal}{RGB}{16,56,60}       % #10383C 主色
\definecolor{myorange}{RGB}{237,109,0}    % #ED6D00 BMEB215
\definecolor{myblue}{RGB}{0,113,162}      % #0071A2 Fig 标签
\definecolor{mygreen}{RGB}{30,130,80}     % 正向结果
\definecolor{myred}{RGB}{180,40,40}       % 负向/注意
\definecolor{mygray}{RGB}{130,130,130}
\definecolor{lightteal}{RGB}{230,240,242} % 浅底色
\definecolor{lightorange}{RGB}{255,245,235}

% ===== 字体 (A0) =====
\renewcommand{\normalsize}{\fontsize{22}{30}\selectfont}
\renewcommand{\large}{\fontsize{28}{36}\selectfont}
\renewcommand{\Large}{\fontsize{34}{42}\selectfont}
\renewcommand{\LARGE}{\fontsize{40}{48}\selectfont}
\renewcommand{\huge}{\fontsize{48}{56}\selectfont}
\renewcommand{\Huge}{\fontsize{56}{64}\selectfont}
\renewcommand{\small}{\fontsize{18}{24}\selectfont}
\renewcommand{\footnotesize}{\fontsize{15}{20}\selectfont}

% ===== 间距 =====
\setlength{\parskip}{0.3em}
\setlength{\parindent}{0pt}
\setlength{\multicolsep}{0.5em}
\raggedbottom

% ===== 节标题 (PPT chevron 风格) =====
\usepackage{titlesec}
\titleformat{\section}[hang]
    {\huge\bfseries\color{myteal}}
    {\thesection}
    {0.8em}
    {}
    [\vspace{-1mm}{\color{myteal}\rule{\columnwidth}{2pt}}]
\titlespacing{\section}{0pt}{0.6em}{0.3em}

\titleformat{\subsection}[hang]
    {\Large\bfseries\color{myteal!85}}
    {\thesubsection}
    {0.8em}
    {}
\titlespacing{\subsection}{0pt}{0.3em}{0.1em}

% ===== 列表 =====
\setlist{nosep, leftmargin=2.5em, itemsep=0.3em}

% ===== 图标签 (PPT: #0071A2) =====
\newcommand{\figcaption}[1]{%
    \vspace{0.3em}
    {\color{myblue}\bfseries #1}
    \vspace{0.2em}
}

% ===== 占位图 =====
\newcommand{\placeholder}[2]{%
    \begin{center}
    \colorbox{mygray!12}{\parbox{0.9\columnwidth}{%
        \centering
        \vspace{1.5em}
        {\Large\bfseries\color{mygray} [ 占位图： #1 ]}
        \vspace{0.8em}
        {\small\color{mygray} #2}
        \vspace{1.5em}
    }}
    \end{center}
}

\begin{document}

% =====================
% 标题区 (PPT 模板风格)
% =====================
\begin{center}
    % 课程信息 (orange accent)
    {\fontsize{24}{30}\selectfont\bfseries\color{myorange}BMEB215}
    {\fontsize{20}{28}\selectfont\color{myteal} \hspace{0.5em} 机器学习与医学工程应用}
    \hspace{2em}
    {\fontsize{18}{24}\selectfont\color{myteal!70} ML\&MEA · 2026 Spring}

    \vspace{12mm}

    % 主标题
    {\fontsize{52}{62}\selectfont\bfseries\color{myteal}
        从任务专用脑电解码器到脑电基座大模型：\\[3mm]
        EEGNet、EEG Conformer 与 CBraMod 的复现与扩展评测
    }

    \vspace{10mm}

    % 作者
    {\fontsize{26}{34}\selectfont\color{myteal!80}
        第一作者 (SUSTech ID)\hspace{2em} 第二作者 (ID)\hspace{2em} 第三作者 (ID)
    }

    \vspace{4mm}

    % 单位
    {\fontsize{20}{28}\selectfont\color{mygray}
        南方科技大学 · 生物医学工程系
    }

    \vspace{8mm}
    {\color{myteal}\hrule height 2.5pt}
    \vspace{6mm}
\end{center}

\begin{multicols}{2}

% ================================================================
% 左栏 · Introduction
% ================================================================
\section{Introduction}

\subsection{为什么需要脑电基座大模型？}

EEG 解码长期面临三个核心困难：

\begin{itemize}
    \item \textbf{标注数据少} —— 多数 BCI、临床、情绪、睡眠数据集被试数量有限，标签成本高。
    \item \textbf{跨域差异大} —— EEG 信号受被试、session、设备、导联、采样率和预处理流程影响明显。
    \item \textbf{评测体系碎片化} —— 不同论文使用不同 split、滤波、输入尺度，模型之间难以公平比较。
\end{itemize}

传统任务专用模型 (CNN/RNN/Transformer) 在特定数据集上从零训练，优势是结构明确、参数量小、单任务性能强；劣势是缺少跨任务复用能力。基座大模型改变的是训练范式：先通过大规模 EEG 数据预训练学习通用表征，再通过 Linear Probing (LP) 或 Fine-Tuning (FT) 适配不同下游任务。

\subsection{发展脉络}

\begin{center}
{\small
\rowcolors{2}{lightteal}{white}
\begin{tabular}{llll}
\toprule
\textbf{阶段} & \textbf{代表工作} & \textbf{为什么选它} & \textbf{主要价值} \\
\midrule
深度学习进入 EEG & EEGNet (2018) & 紧凑、领域感知的深度监督模型 & 高效学习 EEG 特征 \\
全局时序建模 & EEG Conformer (2023) & CNN 前端 + Transformer 注意力 & 长程依赖 + Attention 可视化 \\
大规模预训练 & CBraMod (2025) & EEG Foundation Model & 预训练 → LP/FT 多任务适配 \\
\bottomrule
\end{tabular}
}
\end{center}

\vspace{0.5em}

\textbf{核心问题：}EEGNet 和 EEG Conformer 代表任务专用深度模型，CBraMod 代表从"训练专家模型"到"预训练通用表征再适配下游任务"的范式转变。为什么在专家模型已能在单任务上取得强性能时，我们仍需要脑电基座大模型？

% ================================================================
% 左栏 · Method (NeuroGate 已移出)
% ================================================================
\section{Method}

\subsection{EEGNet —— 深度学习 EEG 解码代表}

用紧凑 CNN 将 EEG 领域知识融入深度学习，是 EEG 解码中最经典的专家模型。

\textbf{架构：}Temporal Convolution (频率时间滤波器) $\rightarrow$ Depthwise Spatial Convolution (通道空间模式) $\rightarrow$ Separable Convolution (降参)。每个下游任务单独监督训练。

\textbf{定位：}任务专用深度模型 baseline，与 EEG Conformer、CBraMod 形成"专家模型 vs 基座模型"对照。

\vspace{0.8em}

\subsection{EEG Conformer —— Transformer 增强专家模型}

在卷积局部特征提取之后引入 Transformer，用 Self-Attention 建模 Patch 间全局依赖。

\textbf{架构：}CNN 前端提取局部 EEG 特征并形成 Patch Embedding $\rightarrow$ Transformer Encoder 建模长程时序依赖 $\rightarrow$ 分类头。Attention Map 可用于可视化分析。

\textbf{定位：}比 EEGNet 更复杂的任务专用 baseline。进行"去掉 Transformer block"消融，检验 Attention 对单任务 EEG 解码的贡献。

\vspace{0.8em}

\subsection{CBraMod —— EEG 基座大模型}

通过大规模预训练获得通用 EEG 表征，再迁移到不同下游任务。

\textbf{下游适配：}CBraMod Backbone 输出 Channel-Patch 表征 $B \times C \times S \times D$。LP 冻结预训练 Backbone，只训练线性分类头；FT 同时更新 Backbone 和分类头。LP 衡量预训练表征本身是否可读，FT 衡量最终下游适配能力。

\vspace{0.8em}

\figcaption{Fig 1. 三类 EEG 解码模型架构对比 (EEGNet / EEG Conformer / CBraMod)。}
\placeholder{模型架构对比示意图}{%
    横向展示 EEGNet (CNN)、EEG Conformer (CNN+Transformer)、\\
    CBraMod (Criss-Cross Attention + Pretraining) 的关键模块和数据流。\quad (等待绘制)}

% ================================================================
% 左栏 · Results
% ================================================================
\section{Results}

\subsection{CBraMod 下游任务能否复现？}

{\footnotesize
\begin{center}
\rowcolors{2}{lightteal}{white}
\begin{tabularx}{\columnwidth}{l>{\raggedright\arraybackslash}p{3.5cm}ccccccc}
\toprule
\multirow{2}{*}{\textbf{数据集}} & \multirow{2}{*}{\textbf{版本配置}} & \multirow{2}{*}{\textbf{Seeds}} &
\multicolumn{2}{c}{\textbf{bAcc}} & \multicolumn{2}{c}{\textbf{\(\kappa\)}} &
\multirow{2}{*}{\textbf{论文参考}} & \multirow{2}{*}{\textbf{结论}} \\
\cmidrule(lr){4-5} \cmidrule(lr){6-7}
& & & Mean & Best & Mean & Best & & \\
\midrule
BCIC2A & session T/E+CAR & 42/2024/3407 & 0.6382 & 0.6709 & 0.5177 & 0.5612 & \(\kappa\) 0.5138 & \cellcolor{mygreen!15}对齐 \\
FACED & baseline official FT & 42/2024/3407 & 0.5245 & 0.5340 & 0.4617 & 0.4718 & bAcc 0.5509 & 接近 \\
ISRUC-S1 & baseline official FT & 42/2024/3407 & 0.7780 & 0.7850 & 0.7128 & 0.7373 & bAcc 0.7865 & \cellcolor{mygreen!15}基本复现 \\
Physionet-MI & flatten FT & 42/2024/3047 & 0.5975 & 0.6181 & 0.4631 & 0.4903 & bAcc 0.6417 & 接近 \\
SEED-V & baseline official FT & 42/2024/3407 & 0.3927 & 0.4018 & 0.2450 & 0.2592 & bAcc 0.4091 & 接近 \\
SHU-MI & flatten FT & 42/2024/3047 & 0.5851 & 0.6380 & 0.1700 & 0.2756 & bAcc 0.6370 & \cellcolor{mygreen!15}best 达到 \\
MDD & flatten FT & 42/2024/3047 & \textbf{0.9034} & \textbf{0.9674} & \textbf{0.8008} & \textbf{0.9099} & bAcc 0.9560 & \cellcolor{mygreen!15}best 超过 \\
TUEV & baseline official FT & 42/2024/3407 & 0.5499 & 0.5755 & 0.5964 & 0.6343 & \(\kappa\) 0.6671 & 未达 \\
\bottomrule
\end{tabularx}
\end{center}
\vspace{0.2em}
{\small\color{mygray}注：bAcc = 平衡准确率，\(\kappa\) = Cohen's Kappa。Mean 为三种子均值，Best 为最优种子。}
}

\figcaption{Fig 2. CBraMod 原论文 8 个下游任务复现全景。多数数据集接近或成功复现论文结果。}

\vspace{1em}

\subsection{EEGNet 复现结果}

{\footnotesize
\begin{center}
\rowcolors{2}{lightteal}{white}
\begin{tabular}{lcccr}
\toprule
\textbf{数据集} & \textbf{类别数} & \textbf{Val Acc (\%)} & \textbf{Test Acc (\%)} & \textbf{Test \(\kappa\)} \\
\midrule
BCICIV2A & 4 & 79.17 & -- & -- \\
BCICIV2B & 2 & 77.72 & 78.35 & 0.5669 \\
IIIa & 4 & 55.95 & -- & -- \\
Physionet MI (22ch) & 4 & 49.30 & -- & -- \\
\bottomrule
\end{tabular}
\end{center}
\vspace{0.2em}
{\small\color{mygray}注：BCICIV2A/BCICIV2B 为主要结果 (300 epochs)；IIIa/Physionet MI 为快速基线 (100 epochs, single seed)。}
}

\vspace{1em}

\subsection{复现中的关键经验}

CBraMod 输出 $B \times C \times S \times D$，其中 $C$ 是通道、$S$ 是 Patch/Segment、$D$ 是隐层维度。如果下游 Head 过早对通道求平均 ($B \times C \times S \times D \rightarrow B \times S \times D$)，会丢失 MI 等任务依赖的空间分布信息。

\begin{center}
{\footnotesize
\rowcolors{2}{lightteal}{white}
\begin{tabular}{lccc}
\toprule
\textbf{数据集} & \textbf{修正前 best bAcc} & \textbf{修正后 best bAcc} & \textbf{论文参考} \\
\midrule
SHU-MI & 0.5499 & 0.6380 & 0.6370 \\
Physionet-MI & 0.5594 & 0.6181 & 0.6417 \\
MDD & 0.9098 & 0.9674 & 0.9560 \\
\bottomrule
\end{tabular}
}
\end{center}

\textbf{可讲结论：}预训练表征不等于下游性能；Foundation Model 还需要与 EEG 结构匹配的 Downstream Readout。

% ================================================================
% 右栏 · Discussion
% ================================================================
\section{Discussion}

\subsection{NeuroGate：重新引入神经科学知识}

随着模型从 EEGNet 到 CBraMod，表征越来越抽象——从可解释的频段特征 (Alpha/Beta/Theta) 变为高维 Foundation Embedding。NeuroGate 尝试将神经科学先验重新注入 Foundation Model。

\textbf{核心思路：}提取人工神经先验特征 $m$（时域统计、Band Power、频谱形状等），通过 Gated Fusion 与 CBraMod 的 Backbone Embedding $h$ 融合：

\[
z = h + g \cdot \phi(m)
\]

目的不是让人工特征替代 Foundation Embedding，而是让模型在需要时利用可解释的 EEG 生理知识。

\subsection{NeuroGate 实验结果}

{\footnotesize
\begin{center}
\rowcolors{2}{lightteal}{white}
\begin{tabular}{lccc}
\toprule
\textbf{数据集} & \textbf{CBraMod LP bAcc} & \textbf{NeuroGate LP bAcc} & \(\mathbf{\Delta}\) \\
\midrule
HMC & 0.6149 \(\pm\) 0.0256 & 0.6922 \(\pm\) 0.0172 & \cellcolor{mygreen!15}\(+\)0.0773 \\
ISRUC-S1 & 0.5962 \(\pm\) 0.0579 & 0.6934 \(\pm\) 0.0156 & \cellcolor{mygreen!15}\(+\)0.0972 \\
MDD & 0.8907 \(\pm\) 0.0477 & 0.9062 \(\pm\) 0.1101 & \(+\)0.0154 \\
Physionet-MI & 0.2719 \(\pm\) 0.0176 & 0.2875 \(\pm\) 0.0083 & \(+\)0.0156 \\
\bottomrule
\end{tabular}
\\[4pt]
{\small\color{mygray}表：冻结骨干网络下的线性探测结果 (LP)。睡眠任务提升最为显著。}
\end{center}
}

{\footnotesize
\begin{center}
\rowcolors{2}{lightteal}{white}
\begin{tabular}{lc}
\toprule
\textbf{实验设置} & \textbf{结果} \\
\midrule
LP channel-gated fusion 平均提升 (vs.\ backbone-only) & \cellcolor{mygreen!15}\(+\)0.0389 BA \\
LP gated fusion 正提升单元比例 (dataset-backbone) & 40/48 \\
CBraMod flatten-FT gated fusion 平均提升 & \(+\)0.0078 BA (9/12 为正) \\
FT best-of-gated/logit 平均提升 & \(+\)0.0097 BA \\
\bottomrule
\end{tabular}
\\[4pt]
{\small\color{mygray}表：Selected-12 数据集上 FT 实验统计结论。}
\end{center}
}

\figcaption{Fig 3. NeuroGate LP 实验结果。显式神经先验能补充 Frozen CBraMod Embedding，在睡眠任务 (HMC, ISRUC-S1) 上提升最显著。}

\subsection{专家模型 vs 基座模型：互补而非替代}

\begin{center}
{\small
\rowcolors{2}{lightteal}{white}
\begin{tabular}{lccc}
\toprule
\textbf{模型} & \textbf{核心思路} & \textbf{泛化能力} & \textbf{可解释性} \\
\midrule
EEGNet & CNN 领域知识 & 低 & 高 \\
EEG Conformer & CNN + Transformer & 中 & 中 \\
CBraMod & 大规模预训练 & 高 & 较低 \\
NeuroGate & + 神经科学先验 & 高 & 预期更高 \\
\bottomrule
\end{tabular}
}
\end{center}

\vspace{0.5em}

\textbf{核心观点：}专家模型衡量单任务上限，基座模型衡量通用表征和迁移能力。两者应互补评价，而非简单替代。EEG Foundation Models 正在解决泛化问题，而可解释性正在成为新的瓶颈——NeuroGate 代表其中一个可能的方向。

\figcaption{Fig 4. 四种 EEG 解码模型架构与范式对比。泛化能力持续提升，可解释性成为新挑战。}

% ================================================================
% 右栏 · Summary
% ================================================================
\section{Summary}

\begin{enumerate}[leftmargin=2em]
    \item \textbf{EEGNet $\rightarrow$ EEG Conformer $\rightarrow$ CBraMod} 代表了 EEG 解码从任务专用专家模型走向基座大模型的范式演进。
    \item \textbf{CBraMod 在多数原论文任务上可复现}，但单任务上限不一定超过精心调参的专家模型——其价值在于跨任务复用和统一表征。
    \item \textbf{下游 Readout 设计至关重要}：保留 Channel-Patch 结构是发挥 Foundation Model 能力的前提。
    \item \textbf{NeuroGate 在 LP 下显著提升}（HMC +0.077, ISRUC-S1 +0.097），为"可解释性增强基座模型"提供了初步证据。
    \item \textbf{统一 Benchmark} 是 EEG Foundation Model 发展的必要基础。
\end{enumerate}

\vspace{0.5em}

\begin{center}
\colorbox{lightteal}{\parbox{0.95\columnwidth}{\centering
    {\LARGE\bfseries\color{myteal}
        EEG Foundation Models 正在解决泛化问题，\\[3mm]
        而可解释性正在成为新的瓶颈。}\\[3mm]
    {\large\color{myteal!80}
        NeuroGate 尝试将神经科学知识重新引入 Foundation Models —— 其中一个可能的方向。}
}}
\end{center}

% ================================================================
% 右栏 · References
% ================================================================
\section*{References}

{\footnotesize
\begin{enumerate}[leftmargin=2em]
    \item Lawhern et al. \textbf{EEGNet: A Compact Convolutional Neural Network for EEG-based Brain-Computer Interfaces}. \textit{J. Neural Eng.}, 2018.
    \item Song et al. \textbf{EEG Conformer: Convolutional Transformer for EEG Decoding and Visualization}. \textit{IEEE TNSRE}, 2023.
    \item Wang et al. \textbf{CBraMod: A Criss-Cross Brain Foundation Model for EEG Decoding}. \textit{ICLR}, 2025.
    \item NeuroGate / NeuroKE 内部参考：neural-prior gated fusion experiments under LP and FT.
\end{enumerate}
}

\end{multicols}

% =====================
% 底部
% =====================
\vspace{4mm}
{\color{myteal}\hrule height 2pt}
\vspace{3mm}
\begin{center}
{\footnotesize\color{mygray}
    BMEB215 · 机器学习与医学工程应用 · 2026 Spring \hspace{3em}
    南方科技大学 · 生物医学工程系
}
\end{center}

\end{document}
